% page 74 du fac-simile (image p-073 du scan)
% bloc 160.0 x 233.3 mm ; marge gauche 8.0 mm
\pgc{-12.700mm}{0.000mm}{
\l{\cel{8}66}
\l{}
\l{\sou{biologio}. Cienco generala pri la vivo, lua kondicioni e lua}
\l{\cel{3}formi diversa (organizeso di la enti vivanta, medii ube}
\l{\cel{3}li vivas, funcioni di la vivo, e c.) - DEFIRS.}
\l{}
\l{\sou{biomekaniko}. Parto di biologio, qua traktas la expliko mekan-}
\l{\cel{3}ista di la fenomeni vivala. - DEFIRS.}
\l{}
\l{\sou{biro}. Drinkajo fermentacinta, preparita ek hordeo ed ek jerm-}
\l{\cel{3}ifinta frumento e lupulo. - DEFI.}
\l{}
\l{\sou{biremo}. Galero qua havas singla-latere du rangi de remisti.}
\l{\cel{3}DEFIRS.}
\l{}
\l{\sou{birko}. (\sou{bot}.)Arboro qua esas un genero di la familio \textquotedbl{}betu-}
\l{\cel{3}lacei\textquotedbl{}. L.\cel{1}\sou{betula}. D.E.}
\l{}
\l{\sou{bireto}. Mikra boneto; kun quar edri quadrata, qua esas fald-}
\l{\cel{3}ebla. - DEFIRS.}
\l{}
\l{\sou{bis}. Duesmafoye. - DFIRS.}
\l{}
\l{\sou{biso}. Materio texebla, sorto di lino flavatra, de qua la An-}
\l{\cel{3}tiqui fabrikis stofi tre granda-preco.- DEFIRS.}
\l{}
\l{\sou{bixestila}.\cel{2}En qua inkluzesis la dio quan onu adjuntas a la}
\l{\cel{3}monato februaro singla-quaresma-yare. - DEFIS.}
\l{}
\l{\sou{bisko}. Kankro-supo. - DEF.}
\l{}
\l{\sou{biskoto}. Loncho de pano preparita per lakto (vice per aquo) e}
\l{\cel{3}sikigita e bakita en forno. - DEF.}
\l{}
\l{\sou{bismuto}. (\sou{kemio}). Metalo grize-blanka kun nuanci redatra,}
\l{\cel{3}lameloza, frajila, qua, amalgamita kun merkurio, uzesas}
\l{\cel{3}por stanizar la speguli. -\cel{1}\sou{Simbolo kemiala}\cel{1}:Bi. - DEFIRS.}
\l{}
\l{\sou{bisquito}. Kukajo lejera, qua konsistas ye ovi, farino e sukro.}
\l{\cel{3}- DEFIRS.}
\l{}
\l{\sou{bistro}. Farbo bruna flavatre, quan onu obtenas de fuligino}
\l{\cel{3}koquita e dilutita, e quan onu uzas por aquo-piktar. -}
\l{\cel{3}DEFR.}
\l{}
\l{\sou{bisturio}. Mikra kultelo kirurgiala, di qua la lamo, movebla,}
\l{\cel{3}mantenesas da resorto kande onu apertas lu. - DEFIRS.}
\l{}
\l{\sou{bito}. (\sou{navig}.)Bloko pozBa en la parto avana di la navo, e qua}
\l{\cel{3}konsistas ye du montanti ed un traverso sur qua spulesis}
\l{\cel{3}e fixigesis per amaro la kabli. - DEFIRS.}
\l{}
\l{\sou{bitra}. Qua saporas repugnante. -(\sou{metaf}.)Qua impresas penigante,}
\l{\cel{3}chagrenigante. - DE.}
}
